\documentclass[12pt,letterpaper]{article}
\usepackage{amsmath}
\usepackage{amsthm}
\usepackage{amsfonts}
\usepackage{amssymb}
\usepackage{amscd}
\usepackage{enumerate}
\usepackage{fancyhdr}
\usepackage{mathrsfs}
\usepackage{bbm}
\usepackage{framed}
\usepackage{mdframed}
\usepackage{cancel}
\usepackage{mathtools}
\usepackage{verbatim}
\usepackage{hyperref}
\usepackage[letterpaper,voffset=-.5in,bmargin=3cm,footskip=1cm]{geometry}
\setlength{\parindent}{0.0in}
\setlength{\parskip}{0.1in}
\allowdisplaybreaks
\headheight 15pt
\headsep 10pt
\newcommand\N{\mathbb N}
\newcommand\Z{\mathbb Z}
\newcommand\R{\mathbb R}
\newcommand\Q{\mathbb Q}
\newcommand\lcm{\operatorname{lcm}}
\newcommand\setbuilder[2]{\ensuremath{\left\{#1\;\middle|\;#2\right\}}}
\newcommand\E{\operatorname{E}}
\newcommand\V{\operatorname{V}}
\newcommand\Pow{\ensuremath{\operatorname{\mathcal{P}}}}

\DeclarePairedDelimiter\ceil{\lceil}{\rceil}
\DeclarePairedDelimiter\floor{\lfloor}{\rfloor}
\newcommand\hint[1]{\textbf{Hint}: #1}
\newcommand\note[1]{\textbf{Note}: #1}
\newcounter{enum22i}
\newenvironment{22enumerate}{\begin{enumerate}[a.]\itemsep0em\setcounter{enumi}{\value{enum22i}}}{\setcounter{enum22i}{\value{enumi}}\end{enumerate}}
\newenvironment{22itemize}{\begin{itemize}\itemsep0em}{\end{itemize}}
\fancypagestyle{firstpagestyle} {
  \renewcommand{\headrulewidth}{0pt}
  \lhead{\textbf{CSCI 0220}}
  \chead{\textbf{Discrete Structures and Probability}}
  \rhead{M. Littman}
}

\fancypagestyle{fancyplain} {
  \renewcommand{\headrulewidth}{0pt}
  \lhead{\textbf{CSCI 0220}}
  \chead{Homework 4}
  \rhead{\textit{ }}
}
\pagestyle{fancyplain}


\begin{document}
  \thispagestyle{firstpagestyle}
  \begin{center}
    {\huge \textbf{Homework 4}}

    {\large Due:\textit{ Monday, June 14th at 11:59pm EDT}}
  \end{center}


     All homeworks are due the Monday after their release at 11:59pm EDT on Gradescope.

    Please do not include any identifying information about yourself in the handin, including your Banner ID.

    Be sure to fully explain your reasoning and show all work for full credit. We recommend checking out the sample proofs on the course website to see the level of rigor for which we're looking.

\subsection*{Problem 1}
{
\nopagebreak 

{\setcounter{enum22i}{0}
\textbf{Note: }For the purposes of this homework, we will also use the relation symbol $R$ to represent the graph of the relation $R$. Thus, $R$ can be used to express relation between two elements, such as $a\;R\;b$, but can also be used to express the graph of the relation, such as $R=\{(1,0),(0,1)\}$.
\begin{22enumerate}

    \item Which of the following relations are symmetric and
    which are antisymmetric? Prove your answers.
    	\begin{enumerate}[i.]
    	\item $R_1 = \setbuilder{(a,b) }{a = b \text{ or } a = -b}$
    	\item $R_2 = \setbuilder{(a,b) }{a = b + 1}$
    	\end{enumerate}
    
    
    \item Consider the following relations on $\mathbb{R}$. State if each relation is a function, a partial function, or not a function, and briefly explain your reasoning. Then:
\begin{22itemize}
    \item If it is a \textit{function}, state if it is injective, surjective, or neither.
    \item If it is a \textit{partial function}, modify the domain to make it a function.  
    \item If it is \textit{not a partial function}, briefly explain why.
\end{22itemize}

\begin{enumerate}[i.]
    \item $R=\{(2,3),(3,4),(4,5),(5,6),(6,7)\}$.
    \item $f(x)=x^2+1$.
    \item $f(x)=y$ if and only if $x^2=y^2$.
    \item $f^{-1}(x)$ where $f(x)=x^5+40$.
     \item $f(x)=\sqrt{x}$.
\end{enumerate}
    
    \item Prove that if a relation $R$ on a set $S$ is reflexive
      and transitive, then the relation $R \cup R^{-1}$ is an equivalence
      relation, where $R^{-1}$ is the inverse.
    
\end{22enumerate}
}
}



\subsection*{Problem 2}
{
\nopagebreak 

{\setcounter{enum22i}{0}

        Let $R$ be an equivalence relation on a finite set $A$, where $|A|=n$ and $n$ is a non-negative integer.
        \begin{enumerate}[a.]
            \item What is the minimum cardinality of $R$, in terms of $n$? Justify (prove) why you can achieve this cardinality and why no smaller cardinality can be achieved.
            \item What is the maximum cardinality of $R$, in terms of $n$? Justify (prove) why you can achieve this cardinality and why no larger cardinality can be achieved.
            \item Suppose $n$ is even. What can you say about the parity of $|R|$? Justify (prove) your response.
            \item Suppose $n$ is odd. What can you say about the parity of $|R|$? Justify (prove) your response.
        \end{enumerate}
}
}
\subsection*{Problem 3}
{
\nopagebreak 

{\setcounter{enum22i}{0}

	\begin{22enumerate}
	\item \begin{enumerate}[i.]
		\item
		Find all relations with domain $\{0,1\}$ and co-domain $\{1\}$.
		\item
		Find all relations with domain $\{1\}$ and co-domain $\{0,1\}$.
		\end{enumerate}
	\item \begin{enumerate}[i.]
		\item
		Find all functions from $\{0,1\}$ to $\{1\}$.
		\item
		Find all functions from $\{1\}$ to $\{0,1\}$.
		\end{enumerate}
	\item Create a bijection between the set $\Z^+$ and the set $\{n\;|\;n=2k \text{ for }k \in \Z^+\}$ whose inputs and outputs are the \textbf{binary representation} of integers.
	\end{22enumerate}
}
}


\subsection*{Mind Bender (Extra Credit)}
{
\nopagebreak

{\setcounter{enum22i}{0}
	
      \begin{22enumerate}

          \item Consider the relation $R$ on a finite set $S$ where $S \subseteq \Z$. Define $R$ such that $\forall a,b \in S$, if $a > b$ or $a < b$ then $(a,b) \in R$. Prove whether or not the complementary relation $\overline{R}$ is reflexive, irreflexive, symmetric, anti-symmetric, or transitive. Assume the complement is taken relative to the universal set $U=S \times S$.\\
          \textbf{Note: }you must include a proof for each property.

        \item 
          Suppose now that the relation $\preceq _{1}$ on set $S$ and relation $\preceq _{2}$ on set $T$ have the same properties as the relation $\overline{R}$ from part (a). If the relation $\preceq$ on the set $S \times T$ is defined by $(s, t) \preceq (u, v)$ if and only if $s \preceq _{1} u$ and $t \preceq _{2} v$ for $s, u \in S$, and $t, v \in T$, show that $\preceq$ also has the same properties as $\overline{R}$.
      \end{22enumerate}
    
}
}

\end{document}